\chapter{Local Simplification}
\label{ch:ch13-local-simplification}

Global analyses find facts, but they never rewrite a single instruction by themselves.
Somewhere a pass must look at \code{x + 0}, \code{y * 1}, \code{0 * f()}, and \code{br true, A, B} and replace each with the cheaper equivalent.
Doing that cleanup late wastes every downstream pass, while doing it locally and repeatedly keeps the whole pipeline honest.
The problem of this chapter is to simplify code within one block, safely and to a fixpoint, before global passes spend budget on it.
We treat each rewrite as a guarded rule, each block as a worklist, and each iteration as progress toward a normal form.
Section~\ref{sec:ch13-rules} catalogs the rules, Section~\ref{sec:ch13-order} orders them, and Section~\ref{sec:ch13-worked} simplifies the block in Figure~\ref{fig:ch13-block} step by step.

\section{Rewrite rules with guards}
\label{sec:ch13-rules}

A \term{rewrite rule} pairs a pattern with a replacement and a guard that states when the replacement preserves behavior.
The pattern \code{x + 0} rewrites to \code{x} with the guard that addition has no trapping overflow on this type.
The pattern \code{x * 0} rewrites to \code{0} with the guard that \code{x} cannot be a trapping value such as infinity, a signaling not-a-number, or a faulting address expression.
The pattern \code{br true, A, B} rewrites to \code{jmp A} unconditionally, because branch folding on a known condition removes no behavior.
The guard is the rule, in the sense that an unguarded rewrite is a bug waiting for its counterexample.
Table~\ref{tab:ch13-rules} lists the core rules used in this chapter with their guards and their payoffs.

Identities remove operations that contribute nothing to the value.
Adding zero, multiplying by one, shifting by zero, masking with all ones, and OR-ing with zero all return their other operand.
Annihilators replace the operation with a constant, for example AND with zero yields zero.
Each identity carries an effect guard: the discarded operand must not trap, fault, or diverge when evaluated.
A pure integer operand passes trivially, a floating operand needs a finiteness or no-trap argument, and a memory operand needs an aliasing and faulting argument.
Beginners drop the guard because tests never trap, and the miscompile appears only on the first signaling input in production.
The C standard makes several of these guards explicit by leaving overflow, division, and aliasing latitude to implementations, see~\cite{wg14-2011-c11}.

\begin{irlisting}
; Block needing several local rules before any global pass runs
; Guards assume 32-bit two's complement ints without trapping overflow
B1: t1 = a + 0
B1: t2 = t1 * 1
B1: t3 = b - b
B1: t4 = t2 + t3
B1: t5 = t4 * 0
B1: c1 = t5 == 0
B1: br c1, B2, B3
\end{irlisting}

The listing compresses five rules and one folding opportunity into seven lines.
The sum with zero, the multiply by one, the self-subtraction, and the multiply by zero each match an algebraic rule.
The comparison of the result against zero then folds to a constant, and the branch folds to a jump.
No fact from outside the block is needed, only the guards on integer purity.
A global pass that ran before this cleanup would analyze seven instructions instead of one jump.
A global pass that runs after sees the normal form and spends its budget on facts only local rules cannot derive.
That division of labor is the economic argument for simplifying early and often.
Algorithm~\ref{alg:ch13-simplify} mechanizes it without rescanning the whole function after each rewrite.

Strength reduction and reassociation sit at the boundary of local power.
Replacing \code{x * 2} with \code{x + x} or \code{x << 1} is local in shape but global in profit, since the target cost depends on the machine.
Reassociating \code{(a + b) + c} into \code{a + (b + c)} to expose a constant fold is likewise tempting and risky.
Floating-point reassociation changes rounding and exception behavior, so the guard must cite the fast-math latitude explicitly rather than assume it.
Integer reassociation is safer but still interacts with overflow wrapping and undefined cases.
Our local rules therefore canonicalize cheaply, for example by sorting constant operands to one side, and leave machine-specific strength choices to later passes.
Morgan places these target-sensitive strength decisions after machine-independent cleanup, see~\cite{morgan-1998-optimizing}.

\section{Ordering rules toward a normal form}
\label{sec:ch13-order}

A \term{normal form} is a shape in which no rule applies, so iteration stops with nothing left to do.
Confluence would mean every rule order reaches the same normal form, which rarely holds once reassociation meets folding.
Termination needs more care: a rule that folds and a rule that unfolds the same pattern can ping-pong forever.
We impose a \term{rule order} with a simple measure that strictly decreases on every rewrite.
Constant folding runs first because it shrinks the program and enables branch folding.
Algebraic identities run next because they remove instructions without creating foldable constants elsewhere in the block.
Canonicalization runs last because it rearranges rather than shrinks, and rerunning folders after it catches newly adjacent constants.
The loop repeats until a full sweep changes nothing, which the measure guarantees will happen.
The measure combines instruction count, constant count, and a canonicalization rank in lexicographic order.
Folding reduces instruction count or replaces an operation with a constant of lower rank.
Identity removal reduces instruction count while never increasing constants elsewhere.
Canonicalization keeps instruction count fixed but lowers rank by moving constants rightward and sorting operands.
No rule increases an earlier component of the triple to decrease a later one, so the triple falls on every accepted rewrite.
A rule that would raise the measure is deferred to a later pass with a different measure, not forced into this loop.
This discipline sounds fussy until the first infinite loop between reassociate-left and reassociate-right teaches it empirically.
Section~\ref{sec:ch13-implementation} shows how the worklist enforces the same discipline incrementally.

\begin{algorithm}
\caption{Local simplification to a normal form within one block}\label{alg:ch13-simplify}
\begin{pseudocode}
\Require block $B$ with instructions $i_1$ to $i_n$, ordered rule set $R$ with decreasing measure
\Ensure $B$ in normal form: no rule in $R$ applies to any instruction
\State $W \gets$ all instructions in $B$
\While{$W \neq \emptyset$}
\State remove some $j$ from $W$
\ForAll{rules $r \in R$ in priority order}
\If{$r$ matches $j$ and its guard holds}
\State rewrite $j$ by $r$
\State add surviving $j$, its users, and new neighbors in $B$ to $W$
\State \Break out of the rule loop to re-examine $j$
\EndIf
\EndFor
\EndWhile
\end{pseudocode}
\end{algorithm}

Algorithm~\ref{alg:ch13-simplify} visits instructions in priority order and requeues only their users and new neighbors.
Rewriting \code{t1 = a + 0} to a copy requeues \code{t2 = t1 * 1} because its operand changed, not the whole block.
Folding a comparison requeues its branch because the branch condition changed from computed to constant.
The \code{break} after a successful rewrite restarts rule matching at highest priority for the same instruction, so folding always precedes canonicalization on fresh patterns.
Termination follows from the decreasing measure on every rewrite plus a finite requeue fan-out per rewrite.
Correctness follows rule by rule from guards, since the loop applies no transformation beyond guarded rules.
The \enquote{local} in the name refers to the scope of facts used, not to the size of the win: branch folding from local constants regularly deletes whole downstream blocks.
Appel implements this cleanup as the peephole layer beneath global optimization, see~\cite{appel-1998-modern}.
\section{Worked example in seven rewrites}
\label{sec:ch13-worked}

We simplify the block from Figure~\ref{fig:ch13-block} by hand under 32-bit wrapping integer guards.
The initial block computes \code{t1} through \code{t5}, compares against zero, and branches.
Table~\ref{tab:ch13-rules} names each rule applied, and the trace below numbers the seven rewrites in order.
First, \code{t1 = a + 0} rewrites to a copy of \code{a}, since integer addition by zero neither traps nor changes the value.
Second, \code{t2 = t1 * 1} rewrites to a copy of \code{a} after copy propagation of the new \code{t1}, since multiply by one is an identity on integers.
Third, \code{t3 = b - b} rewrites to constant zero, since self-subtraction of an integer is zero without reading memory.
Fourth, \code{t4 = t2 + t3} becomes \code{a + 0} after substituting the copies and the zero, then rewrites to \code{a} by the additive identity.
Fifth, \code{t5 = t4 * 0} rewrites to constant zero, with the pure-integer guard discharging the annihilator safely.
Sixth, \code{c1 = t5 == 0} folds to constant true, since both operands are now known.
Seventh, \code{br c1, B2, B3} folds to \code{jmp B2}, deleting the edge to $B3$ and exposing it to dead-code removal.
The block that began as seven instructions ends as a direct jump with all temporaries dead.

\begin{figure}
\begin{tikzpicture}
\node[basic block] (b0) {entry: a, b live};
\node[basic block, below=1.0cm of b0] (b1) {B1: t1..t5, cmp, br};
\node[basic block, below left=1.0cm and 0.6cm of b1] (b2) {B2: use a};
\node[basic block, below right=1.0cm and 0.6cm of b1] (b3) {B3: use b};
\draw[cfg edge] (b0) -- (b1);
\draw[cfg edge] (b1) -- (b2);
\draw[cfg edge] (b1) -- (b3);
\end{tikzpicture}
\caption{The single simplifiable block $B1$ with two successors, where folding the branch to $B2$ later deadens the edge to $B3$.}
\label{fig:ch13-block}
\end{figure}

\begin{table}
\caption{Local rules applied in the worked example, each with its guard for wrapping integer arithmetic.}
\label{tab:ch13-rules}
\begin{tabular}{lcc}
\toprule
Pattern and replacement & Guard & Effect \\
\midrule
x + 0 becomes x & no trapping overflow & removes add \\
x times 1 becomes x & integer type & removes mul \\
x - x becomes 0 & pure integer operand & creates constant \\
x times 0 becomes 0 & pure, no trap & creates constant \\
cmp folded, br folded & constant condition & removes edge \\
\bottomrule
\end{tabular}
\end{table}

Each step preserves the observable behavior while strictly lowering the measure.
The first four steps lower instruction count by removing arithmetic.
The fifth replaces a multiply with a constant of lower rank.
The sixth lowers rank by turning a comparison into a constant.
The seventh lowers count by replacing a conditional branch with a jump.
No step uses a fact from $B2$ or $B3$, and no step needs the global analyses from earlier chapters.
A checker can verify the sequence by replaying the seven rewrites and confirming the guard at each one.
The most instructive check is the fifth: replacing it with a floating multiply would demand a finiteness guard the example does not supply.
That failed guard is exactly where fast-math latitude would have to appear explicitly, and its absence keeps the integer rule honest.

The example shows why simplification is useful before global passes; running it again afterward can exploit facts those passes discover.
Available-expression analysis on the original block would number five temporaries and two comparisons.
Liveness would keep \code{a} and \code{b} live through all seven lines.
After simplification, both analyses see one jump and nothing else.
Compile time saved downstream usually exceeds the linear cost of the local sweep.
More importantly, the answers improve: constants discovered locally seed the global constant propagator of the next chapter.
Safonov treats this layering as part of predictable compiler architecture rather than as opportunistic cleanup, see~\cite{safonov-2010-trustworthy}.

\section{Implementation notes and cost}
\label{sec:ch13-implementation}

An implementation stores one rule table shared by all blocks, with pattern matchers compiled to a decision tree over opcodes.
Matching tries folders first, then identities and annihilators, then canonicalization, exactly mirroring the priority in Algorithm~\ref{alg:ch13-simplify}.
Guards consult type tags, fast-math flags, and a purity bit on each operand, all available locally.
Rewrites update def-use chains in place and requeue users, so the per-rewrite cost is proportional to the user count.
Let $n$ be the initial instruction count, $r$ the successful rewrite count, and $u$ the maximum number of affected users and neighbors per rewrite.
With constant-cost rule matching, duplicate-free queue membership, and constant-cost bookkeeping per affected instruction, the cost is $O(n+r(u+1))$.
A well-founded decreasing measure proves termination but does not by itself bound $r$ by $O(n)$; that stronger bound needs a rule-specific argument.
Memory overhead is one worklist plus def-use edges the compiler already keeps.

The main pitfall is rule interaction across passes, not within one block.
A canonicalizer that sorts operands can fight a selector that wants constants on a specific side for addressing modes.
A folder that creates \code{0.0} with the wrong sign can break a later sign-sensitive branch fold.
A demolished type guard can let integer rules fire on pointers, miscompiling address arithmetic the C standard reserves to implementations.
The fix layers measures: local simplification owns the portable measure, while target passes own their own measures afterward.
Rules that need target costs carry explicit hooks rather than smuggling prices into guards.
Testing covers each rule with pure, trapping, floating, and effectful operands so guards fire exactly when they should.
Section~\ref{sec:ch13-rules} lists the guard for every rule, and the exercises ask you to break each one deliberately.
% TODO(cite): need a primary reference for peephole and rule-based simplification catalogs beyond the four listed books.

\section{Exercises}
\label{sec:ch13-exercises}

\begin{enumerate}
\item \label{ex:ch13-guard-break} Revisit Table~\ref{tab:ch13-rules} and construct a counterexample for each rule with its guard removed.
Use a signaling not-a-number for the floating annihilator, a faulting load for the identity through memory, and a wrapping versus trapping type pair for the additive identity.
Show the original and rewritten behavior on your counterexample and state which observable event differs.
Then restore the guard and explain why it excludes your case.
\item \label{ex:ch13-measure-fight} Implement Algorithm~\ref{alg:ch13-simplify} with reassociate-left and reassociate-right both enabled at equal priority.
Run it on Figure~\ref{fig:ch13-block} extended with \code{t6 = (a + b) + c} and count rewrites to termination or loop.
Then impose the priority from Section~\ref{sec:ch13-order} and report the new count with the final normal form.
Explain how the lexicographic measure rules out the ping-pong.
\item \label{ex:ch13-branch-cascade} Extend the worked example so $B2$ begins with \code{d = a + 0} and ends with a foldable branch.
Simplify $B1$ to a jump, delete the edge to $B3$, run dead-code reasoning informally, then simplify $B2$ with the newly constant context.
Report how many instructions the cascade removes beyond the seven local rewrites.
Discuss where a global worklist should resume after a local fold changes the CFG.
\end{enumerate}

\section{Further reading}
\label{sec:ch13-reading}

Appel treats local cleanup as the peephole foundation under global passes, see~\cite{appel-1998-modern}.
Morgan separates machine-independent simplification from target-sensitive strength choices, see~\cite{morgan-1998-optimizing}.
Safonov describes how layered passes keep a trusted pipeline predictable, see~\cite{safonov-2010-trustworthy}.
The C language definition supplies the overflow, rounding, and aliasing rules that guards must cite, see~\cite{wg14-2011-c11}.
